\subsection{Machine learning for public budgets and expenditure}
Most work forecasts \emph{aggregate} public spending as a time series. G\"ur et al.~\cite{gur2025} compare machine- and deep-learning forecasters on T\"urkiye's monthly budget expenditure (2008--2024), including crisis periods; Dr\u{a}gulin et al.~\cite{dragulin2025} benchmark LSTM and TCN networks (among other models) against SARIMA, VAR and ECM; and Yang et al.~\cite{yang2023} forecast the government expenditure of 15 large economies with a GRU network and find it superior to ARIMA, ETS, SVR and XGBoost baselines. Barnes et al.~\cite{barnes2023} study empirically how OECD governments reallocate spending across COFOG functions (1996--2017) and find that countries with sounder governance and tighter fiscal rules reallocate more, but they build no predictive model. Nzazi et al.~\cite{nzazi2025} provide a narrative review of AI for public expenditure analysis in low-income countries and conclude that applications are mostly conceptual or pilot-scale. Santschi et al.~\cite{santschi2024} train XGBoost and CatBoost on Swiss cantonal accounts and study human--AI collaboration against expert forecasters (a short paper with preliminary results); the setting is one country and account groups rather than functional classes. R\u{a}dulescu et al.~\cite{radulescu2015} forecast public expenditure by function with feed-forward networks for five Central and Eastern European economies, before gradient-boosted trees became standard. Valle-Cruz et al.~\cite{valle2022} combine neural networks and a multi-objective genetic algorithm on 217 countries to relate expenditure strategies to outcomes such as inequality and growth and to explore how resources could be allocated; they do not predict an observed functional split.

\subsection{Funding allocation and public spending outcomes}
Caravaggio et al.~\cite{caravaggio2025} predict which actors win competitive EU rural-development funds; Di Stefano and Resce~\cite{distefano2025} uncover a need--allocation paradox in Italian recovery-plan calls; and Chand and Zhang~\cite{chand2022} study the gap between funding and expenditure in Australia's disability insurance scheme. These studies predict \emph{who} receives money at project or participant level, whereas we predict \emph{how a government budget is divided across functions}. On the outcome side, Guerrero and Casta\~neda~\cite{guerrero2022} propose a bottom-up causal framework linking spending to development gaps in 140 countries, and Nguyen et al.~\cite{nguyen2026} document a non-linear association between government expenditure and SDG progress. We use the SDGs only as a \emph{reporting} cross-walk from COFOG functions (Table~\ref{tab:sdg}; cf.\ the 2030 Agenda~\cite{un2015agenda}).

\subsection{Determinants and dynamics of public spending}
The public-economics literature studies why spending levels and mixes differ across countries and years. Government size has been analysed in the tradition of Wagner's law~\cite{lamartina2011}; the composition of public expenditure has been related to economic growth~\cite{devarajan1996,kneller1999,afonso2010}; globalisation, military-spending determinants and budget institutions have been examined as drivers of spending patterns~\cite{garrett2001,dunne2003,alesina1999}; and ideology and electoral cycles have been linked to the \emph{composition} of spending~\cite{potrafke2010,bove2017,nordhaus1975,rogoff1988,brender2005}. Budgeting has long been described as incremental~\cite{davis1966}, with occasional large punctuations~\cite{jones2005}; persistence as a hard-to-beat baseline is therefore the \emph{expected} result of this literature, and we test it rather than present it as a discovery. These studies are explanatory and use regressions on aggregates or a few categories; they do not provide out-of-sample prediction of the full functional composition. We add fiscal-rule, election and IMF-programme variables in one robustness module (Section~\ref{sec:ext}).

\subsection{Methodological background}
\textbf{Compositional targets.} Shares of a total are compositional data~\cite{aitchison1986}; log-ratio transformations~\cite{egozcue2003} and Dirichlet or beta regression~\cite{ferrari2004,douma2019} are the standard tools, and compositional methods have been applied to budget categories~\cite{adolph2020}. \textbf{ML in economics and forecasting.} Overviews of ML for economists~\cite{athey2019,mullainathan2017} and studies of ML for macroeconomic forecasting~\cite{medeiros2021,goulet2022} report that nonlinear learners can improve on standard benchmarks when data are rich, which motivates, but does not guarantee, gains on budget shares. \textbf{Trees versus deep learning on tabular data.} Tree ensembles remain strong on medium-sized tabular data~\cite{grinsztajn2022,shwartzziv2022,borisov2024}, and tabular deep models need careful design~\cite{gorishniy2021}; our panels are small, so we include an MLP and a feature-tokenising transformer only as comparators. \textbf{Uncertainty.} Split-conformal prediction gives distribution-free intervals under exchangeability~\cite{lei2018,angelopoulos2023}, and variants adapt the interval width to heteroscedasticity~\cite{romano2019}; coverage can fail when exchangeability fails~\cite{barber2023}, as it does when a new economy differs from the calibration economies. \textbf{Fairness.} Error that differs systematically between groups is an equity issue~\cite{mehrabi2021}; deployed prediction systems have been shown to under-serve disadvantaged groups~\cite{obermeyer2019}, and error-rate parity is one formal criterion~\cite{hardt2016}.

\subsection{Explainable AI in fiscal applications}
SHAP~\cite{lundberg2017,lundberg2020} and LIME~\cite{ribeiro2016} are routinely attached to gradient-boosted models in economics and public-sector studies, for instance for enterprise performance~\cite{zhang2026} or for patient-level health-insurance cost prediction~\cite{yao2025}. Such studies usually display attribution plots without testing whether the attributions are \emph{faithful, stable or consistent across methods}, although explainers are known to disagree with each other~\cite{krishna2022}, post-hoc explanations can be unstable~\cite{alvarezmelis2018,ghorbani2019}, manipulable~\cite{slack2020} and, in high-stakes settings, inferior to inherently interpretable models~\cite{rudin2019}; Shapley-value attributions have specific conceptual problems as importance measures~\cite{kumar2020}; and removal-based evaluations of attributions, which our deletion test resembles, have their own pitfalls~\cite{hooker2018,bhatt2020}. We follow the call for rigorous evaluation of interpretability~\cite{doshi2017} and report agreement, stability, deletion and ground-truth tests.

\subsection{Recent closely related studies (2025--2026)}
A web search in November 2026 (strings in the repository; a Scopus, Web of Science and IEEE Xplore search is still to be run) found five recent studies that touch our setting. Nonato et al.~\cite{nonato2026} predict Brazil's federal executed spending by budgetary function with a multi-output XGBoost and Bayesian optimisation; their series has 24 annual observations that they augment with 1\,048 synthetic points, they report $R^2$ up to 0.97 and a separate leakage-free check with much lower performance (we could verify the first figure but not the second from the accessible record). The study covers one country and does not use COFOG; as far as we can tell from the abstract it uses no explanation method. Kisman et al.~\cite{kisman2026} explain cross-country defence expenditure (2002--2023) with explainable-AI methods; it targets one function and an expenditure level rather than the composition of the whole budget. Dhungel et al.~\cite{dhungel2025} report an exploratory case study of AI decision support for budget preparation in Germany, and Francies et al.~\cite{francies2026} propose a SHAP, causal and counterfactual framework for urban budget allocation in smart cities (preprint). Mohith et al.~\cite{mohith2026}, known to us from the abstract only, benchmark classical and deep-learning models for one-year-ahead forecasts of India's Union-Budget allocations to twelve sectors (FY 2005--06 to 2024--25), find that gradient boosting and ridge regression beat the deep-learning models, and add SHAP, permutation importance, LIME and prediction intervals. Their finding that simple models win on short annual series agrees with ours, as does the caution about leakage in Nonato et al. What we add is scope and protocol: an \emph{official multi-country} COFOG panel with leave-country-out and temporal splits, persistence and peer baselines reported alongside every learner, a cluster-bootstrap and Holm-corrected comparison, and a test of the explanations rather than their display.

\subsection{Positioning}
We do not claim to be the first to predict budget allocations by ML, nor to propose a new learning method. Our contribution is the \emph{conjunction} of (i) an official, cross-validated multi-country functional-composition target, (ii) practitioner baselines and corrected inference, (iii) quantitative evaluation of SHAP and LIME on it, and (iv) an open audit of a leaky precursor benchmark. To the best of our (non-systematic) search no single prior study covers all four, but each is individually covered by some of the work above. Table~\ref{tab:gap} lists objective attributes of the closest studies; cells are filled from the full text where we read it and otherwise marked as judged from the abstract.
